%! Author = jacopo
%! Date = 2/25/26

% Preamble


\section{Related Work}

\subsection{EEG Foundation models}
Foundation models are machine learning models trained on vast, immerse datasets.
The ones that adopt contrastive or self supervised learning techniques acts as building blocks for specified tasks
for the general latent space they map to.
There have been various proposals for EEG foundation models.
These, most times, suffer from limitations that do not necessarly derive from their architecture but are directly linked
to the size and task specifity of their datasets. The measurements themselves for EEG data can vary much based on
the technology used as there are some standards but they differ between themselves.
The measurments layout used for AMIGOS with the (aparatus name) does share some nodes with the one used in DEAP
(nome di questo) but they mostly track different electrodes.
Thus the cross-dataset compatibility tends to be lower than for other common tasks like vision and audio.

\texit{CBraMod} has been designed to tackle this very problem, its representations are in fact not tied to the single
electrode but to the relative space the nodes occupy.
It tries to run on a higher level and try to do brain-region level modelling.

Some limitations are hard to tackle without an obvious solution: more data.
A persons \textit{EEG} is very personal. It can differ much from one of another's person.
Unless a novel approach is able to abstract the personal encoding inside of the data most solutions will encode the
identity features in the learnt representations.
This is the reason for, in the field, evaluations are more commonly ran subject wise and not experiment wise.

One more thing to consider is that EEG data measurements tends to be noisy and are very artifact prone.
So, despite advances regarding spatial modelling, current EEG foundation models are still limited by dataset scale,
inter-subject variability and predominant unimodal training paradigms.

\subsection{Multimodal Contrastive Models}
Gathering more data is always a costy task and not always sustainable.
In the field of electroencelography even more so.
So the idea to draw more information from different data sources could increase the data that is available to train a model.

Such apporaches have been widely adopted in many fields.
\texttt{Flamingo} by Google Deepmind is an application that partially inspired this work.
It is a Visual-Language model that uses two frozen backbones for image and text generation.
It leverages information given by the visual encoder to enrich the textual representations.
The model requires the two data sources to be interleaved to encode the spatio-temporal information that derives
from the visual encoder inside the correct part of text.
The visual tokens generated by the corresponding backbone are used to condition the text with \textit{cross attention} mechanism.
This is applied in an interleaved way between the layers defined inside of the frozen LM and the new atchitecture.

This way the model is able to preseve the ability it gained in the pretrained LM while injecting the multimodal information
progressively. This also limits the destruction on the well structure data that derives from the originally frozen LM backbone.

A similar commonly considered safer bet is CLIP stlye models which resembles more our resulting model.
Where instead of fusing the information in the whole architecture it is performed only at the end.
It is safer in the regard of overfitting as data limitations or noise in \textit{Flamingo} can have catastrophic
consqeunces to the model general performance.
CLIP also solves a problem posed by \textit{Flamingo} in our domain, it is completely unsupervised.
While the streams of data increase in cases of unsupervised learning CLIP style losses train the model on the batch strcture.
-definizione di clip stlye losses come InfoNCE-.

-- comprimi MOCO --
To solve limitations posed by CLIP style losses and contrastive learning the -ente- came up with \textit{Momentum Contrast (MOCO)}.
It builds a dynamic dictionary with a queue and a moving-averaged encoder.
It has the ability to reduce the memory load of the model during training by generating a queue of computed samples to
run the contrastive loss on.
This way, when in constrained environments, the \textit{batch size} limits can be mitigated by having a representaitonal
proxy of the past processed samples. It virtually enabled to compute the loss on a larger domain than possible without
having to increase the batch size.
To avoid the encodings to be too sparse and volatile the introduced framework makes use of a second instance of the model
that moves in the same direction as the learned target model. The moment is moderated by a \textit{momentum} parameter
that decides the rate of change of the moving model. This allows to tweak the queue to be effective w.r.t. the training objective.

While multimodal contrastive frameworks have affirmed themselves in many prominent fields demonstrating strong cross modal
alignment, their application on EEG representation learning remain limited, particularly under data scarce conditions.

\subsection{Knowledge Distillation}
Another novel technique that aims to improve the training of a model is \textit{Knowledge Distillation (KD)}.
The idea of \textit{KD} is simple: deploy a teacher in a process where its learnt representations are transferred to a target student.
Its most common use case is when compressing architectures.
Following the logic of dead networks most neural models tend to have many dead weights.
It is also often needed to compress the models to be less demanding in terms of resources, in fact many large models
have distilled variants.

Distillation as a framework can happen in different points of a model's architecture.
It can also happen in recursive fashion as self-distillation.

In the case the two matching modalities do not belong to the same modality it is considered \textit{cross modal distillation}.
Beyond any nouance and different categorization KD serves as another branch of computation for a model in order to control the signal of backpropagation during training.
This way the information contained in the teacher is transferred indirectly to the student.
The student model is trained to approximate the teacher’s representations, thereby transferring higher-level structural
information into a more constrained architecture.
For the distillation to work it is necessary that the passing of both models work on the same set of data, or else alignment
would be pointless.
VATE, a multimodal dataset for video, audio and text in affective computing domain provides a suitable teacher source
for cross-modal distillation into EEG representations.

The technique in cross-modal applications is being explored more and more in many contexts.
This is not yet so for EEG representation learning even if the data is most times available.
For this reason I consder that investing time in the process could lead to further advances.


Taken togheter the integration of these techniques remain unsuficciently explored.